\begin{pastexam}[label={px:ntu-115-bmei-alg-01}]{臺大生醫電資所 115 學年度 第 1 題（節錄第 (1)(5)(6)(7)(11) 小題）}
（名詞解釋；每小題 2 分。 其餘十五個名詞是計算機概論，未收）

(40\%, 2 points for each problem) Please define the following terms and explain the content,
purpose, and application of each term and give an illustrative example if possible. If possible,
define the term in mathematical equation. If it is an acronym, please write the full name. For
example: SSD: Solid-State Drive: A solid-state drive (SSD) is a storage device that typically uses
flash memory to store data, instructions, and information. Flash memory contains no moving parts,
making it more durable and shock resistant than other types of media. SSDs offer higher
input/output rates and lower latency, but are slightly more expensive for the same capacity. For
this reason, manufacturers typically offer SSDs as an option instead of hard disks in their
laptops, tablets, and desktops.

(1) Viterbi algorithm \quad (5) min-max heap \quad (6) AVL tree (Adelson-Velsky and Landis tree)
\quad (7) dynamic programming \quad (11) stack data structure
\end{pastexam}

\begin{pastsoln}
五個小題分屬第~18~章（(1)(7)）、第~8~章（(5)）、第~10~章（(6)）、第~6~章（(11)），放在動態規劃這
一章。 題目要求「定義、內容、用途、應用、例子，能寫數學式就寫」，每個名詞各寫一段。

\textbf{(1) Viterbi algorithm（維特比演算法）}。 在隱藏馬可夫模型（HMM）裡，給定觀測序列
$o_1, o_2, \ldots, o_T$，找出機率最大的隱藏狀態序列的動態規劃演算法。 模型有 $N$ 個狀態、初始機率
$\pi_i$、轉移機率 $a_{ij}$、發射機率 $b_j(o)$。 令 $\delta_t(j)$ 是「前 $t$ 個觀測、第 $t$ 步停在狀態 $j$」
的所有路徑裡最大的機率：
\[
\delta_1(j) = \pi_j\, b_j(o_1), \qquad
\delta_t(j) = \max_{1 \le i \le N} \bigl(\delta_{t-1}(i)\, a_{ij}\bigr)\, b_j(o_t) \quad (t \ge 2),
\]
同時記下取到最大值的 $i$（回溯指標 $\psi_t(j)$）。 最後取 $\max_j \delta_T(j)$，順著 $\psi$ 倒回去就是最可能
的狀態序列。 時間 $O(TN^2)$，若暴力列舉 $N^T$ 條路徑則是指數時間。 應用：語音辨識、詞性標註、
通訊的迴旋碼解碼、生物序列分析（基因預測）。 例子：觀測到病人連續三天「正常、發冷、頭暈」，
推算三天各是「健康」或「發燒」最可能的組合。

\textbf{(5) min-max heap（最小最大堆積）}。 一種雙端優先佇列（資結書第~8.5~節、定義~8.2）：外形是
完整二元樹，各層輪流是 min 層和 max 層：根所在的深度 $0$ 是 min 層，深度 $1$ 是 max 層，依此類推；min 層的節點不大於它子樹裡所有的鍵，
max 層的節點不小於它子樹裡所有的鍵。 所以根是最小值，最大值是根的兩個子節點中較大的那個。
取最小值、取最大值是 $O(1)$，插入、刪除最小值、刪除最大值都是 $O(\lg n)$。 用途：同時需要最小和
最大的場合，例如保留前 $k$ 名又要隨時淘汰最差者的排行榜、外部排序時的緩衝區。 例子：
$\{5, 70, 40, 30, 9, 7, 45\}$ 可以排成根 $5$；深度 $1$ 是 $70$、$45$；深度 $2$ 的 $30, 9$ 在 $70$ 底下，$7, 40$ 在
$45$ 底下。 根 $5$ 是最小，最大值是 $\max(70, 45) = 70$。

\textbf{(6) AVL tree（Adelson-Velsky and Landis tree）}。 一種自我平衡的二元搜尋樹（資結書第~10.2~節、
定義~10.1）：每個節點左右子樹的高度差（平衡因子 $\mathit{bf}(x) = h(x.\mathit{left}) - h(x.\mathit{right})$）
只能是 $-1$、$0$、$1$。 插入或刪除後若某節點的平衡因子變成 $\pm 2$，用 LL、RR、LR、RL 四種旋轉修
回來。 $n$ 個節點的 AVL 樹高度（層數）小於 $1.4405 \lg(n + 1)$（資結書定理~10.1），所以搜尋、插入、刪除都是 $O(\lg n)$，不會像
普通 BST 退化成鏈。 應用：查詢多、需要保證最壞時間的字典或索引，例如記憶體內資料庫索引。 例子：
依序插入 $1, 2, 3$，普通 BST 是一條鏈；AVL 樹在插入 $3$ 時根的平衡因子變 $-2$（RR 型），左旋一次
變成根 $2$、左 $1$、右 $3$。

\textbf{(7) dynamic programming（動態規劃）}。 把問題拆成重疊的子問題，每個子問題只解一次並把答案
存表，再由小問題的答案組出大問題的答案（資結書第~18.1~節）。 適用條件有兩個：最佳子結構（最佳解
由子問題的最佳解組成）和重疊子問題（同一個子問題會被重複用到）。 寫法有由上而下的記憶化和由下而上
填表兩種。 例子：費氏數 $F(n) = F(n-1) + F(n-2)$，直接遞迴是指數時間，填表 $F(0), F(1), \ldots,
F(n)$ 只要 $O(n)$；最長共同子序列
\[
c[i][j] = \begin{cases} 0 & i = 0 \text{ 或 } j = 0, \\ c[i-1][j-1] + 1 & x_i = y_j, \\
\max(c[i-1][j],\, c[i][j-1]) & x_i \ne y_j, \end{cases}
\]
時間 $O(mn)$。 應用：序列比對（生物資訊的 DNA 比對）、背包問題、最短路徑（Floyd--Warshall），
以及上面的 Viterbi 演算法。

\textbf{(11) stack data structure（堆疊）}。 只能在同一端（頂端）放入和取出的線性串列，後進先出
（LIFO）（資結書第~6.1~節）。 基本操作是 push（放到頂端）、pop（取出頂端）、top（看頂端），用陣列或
鏈結串列做都是 $O(1)$。 應用：函式呼叫（呼叫堆疊存返回位址和區域變數）、運算式求值與中序轉後序、
括號配對、DFS 與回溯、瀏覽器的「上一頁」、編輯器的復原。 例子：依序 push $1, 2, 3$，再 pop 兩次，
得到 $3, 2$，堆疊裡剩 $1$。
\end{pastsoln}
